% !TeX root = ../sections/02_exercises.tex
\subsection{1-B-3.}
\begin{exercise}
	以下の問いに答えよ。

	\begin{enumerate}
		\item $n\in\mathbb{N}$ に対して，
		\[
			h_n:=n^{1/n}-1
		\]
		とするとき，$n\geq 2$ に対して不等式
		\[
			h_n^2\leq \frac{2}{n-1}
		\]
		を証明せよ。

		\item
		\[
			\lim_{n\to\infty} n^{1/n}=1
		\]
		が成立することを $\epsilon N$ 論法により証明せよ。
	\end{enumerate}
\end{exercise}
\begin{background}
	この問題では，
	\[
		n^{1/n}\to 1
	\]
	を示すために，
	\[
		h_n=n^{1/n}-1
	\]
	とおいて，$h_n$ が $0$ に近づくことを示す。

	\begin{remark}
		$n\geq 1$ のとき
		\[
			n^{1/n}\geq 1
		\]
		であるから，
		\[
			h_n=n^{1/n}-1\geq 0
		\]
		である。
	\end{remark}

	\begin{theorem}[二項定理]
		$n\in\mathbb{N}$ に対して
		\[
			(1+x)^n
			=
			\sum_{k=0}^{n}\binom{n}{k}x^k
		\]
		が成り立つ。
	\end{theorem}

	\begin{remark}
		$h_n=n^{1/n}-1$ なので，
		\[
			n^{1/n}=1+h_n
		\]
		である。
		したがって
		\[
			n=(1+h_n)^n
		\]
		と書ける。
	\end{remark}
\end{background}
\begin{strategy}
	まず
	\[
		h_n=n^{1/n}-1
	\]
	とおく。

	このとき
	\[
		n=(1+h_n)^n
	\]
	である。

	二項定理を用いると，
	\[
		(1+h_n)^n
		=
		1+nh_n+\frac{n(n-1)}{2}h_n^2+\cdots
	\]
	となる。

	ここで $h_n\geq 0$ であるから，各項は非負である。
	したがって特に
	\[
		n=(1+h_n)^n
		\geq
		\frac{n(n-1)}{2}h_n^2
	\]
	が成り立つ。

	両辺を整理すると，
	\[
		h_n^2\leq \frac{2}{n-1}
	\]
	を得る。

	次に，この不等式から
	\[
		0\leq h_n\leq \sqrt{\frac{2}{n-1}}
	\]
	である。

	右辺は $n\to\infty$ で $0$ に近づく。
	したがって，任意の $\epsilon>0$ に対して，
	\[
		\sqrt{\frac{2}{n-1}}<\epsilon
	\]
	となるように $N$ を選べば，
	$n\geq N$ に対して
	\[
		|n^{1/n}-1|=h_n<\epsilon
	\]
	が従う。

	よって
	\[
		\lim_{n\to\infty}n^{1/n}=1
	\]
	である。
\end{strategy}